GLWE 密码体制是同时涵盖 LWE 与 RLWE 的一般化形式。GLWE 密文是元组 $(\{A_i\}_{i=0}^{k-1}, B)$，其中 $B = \sum\limits_{i=0}^{k-1}{(A_i \cdot S_i)} + \Delta \cdot M + E$。公开量 $\{A_i\}_{i=0}^{k-1}$ 和秘密密钥 $\{S_i\}_{i=0}^{k-1}$ 都是由 $k$ 个 $(n-1)$ 次多项式组成的列表；消息 $M$ 与噪声 $E$ 各为一个 $(n-1)$ 次多项式。与 LWE 和 RLWE 类似，每个密文都使用新的公开随机掩码 $A$，而所有密文使用相同的秘密密钥 $S$。为简化记号，本节将每个密文实例记为 $(\{A_i\}_{i=0}^{k-1}, B)$，而不使用 $(\{A_i\}_{i=0}^{k-1, \langle j \rangle}, B^{\langle j \rangle})$。




\subsection{参数设置}
设 $t$ 为明文空间大小，$q$ 为密文空间大小，其中 $t < q$，即 $t$ 远小于 $q$，且 $t | q$，即 $t$ 整除 $q$。随机选取由 $k$ 个 $(n-1)$ 次多项式构成的列表作为秘密密钥，其中每个系数都是从 $\{-1, 0, 1\}$ 中随机选取的三元数，即 $\{S_i\}_{i=0}^{k-1} \xleftarrow{\$} \mathcal{R}_{\langle n, \textit{tern} \rangle}^k$。令明文尺度因子为 $\Delta = \left\lfloor\dfrac{q}{t}\right\rfloor$。

GLWE 的设置参数与 RLWE 类似。区别在于，$S$ 不再是包含 $n$ 个秘密系数的单个 $(n-1)$ 次多项式，而是由 $k$ 个这样的 $(n-1)$ 次多项式组成的列表，总共包含 $n \cdot k$ 个秘密系数。


\subsection{加密}
\label{subsec:glwe-enc}

多项式代表元的次数至多为 $n-1$，允许高次系数为零。噪声系数是采样后模 $q$ 约减的整数，采用 \autoref{subsec:lwe-dec} 中的约定。正确性不等式约束这些提升后的系数，而不是对无界高斯分布给出确定性保证。

假设有一个 $(n-1)$ 次多项式 $M \in \mathcal{R}_{\langle n, t \rangle}$，其系数表示要加密的明文数值。


\begin{enumerate}
\item 随机选取由 $k$ 个 $(n-1)$ 次多项式组成的列表 $\{A_i\}_{i=0}^{k-1} \xleftarrow{\$} \mathcal{R}_{\langle n,q \rangle}^{k}$，作为一次性的公开随机量。
\item 随机选取小系数多项式 $E \xleftarrow{\chi_\sigma} \mathcal{R}_{\langle n,q \rangle}$ 作为一次性噪声，其 $n$ 个系数都是 $\mathbb{Z}_q$ 中从高斯分布 $\chi_\sigma$ 随机采样的小数值。
\item 将 $M$ 按 $\Delta$ 缩放，即计算 $\Delta \cdot M$。这将 $M \in \mathcal{R}_{\langle n, t \rangle}$ 转换为 $\Delta \cdot M \in \mathcal{R}_{\langle n,q \rangle}$。
\item 计算 $B = \sum\limits_{i=0}^{k-1}{(A_i \cdot S_i)} + \Delta \cdot M + E \in \mathcal{R}_{\langle n,q \rangle}$。
\end{enumerate}

$ $

GLWE 加密公式汇总如下：

$ $

\begin{tcolorbox}[title={\textbf{\tboxlabel{\ref*{subsec:glwe-enc}} GLWE 加密}}]

\textbf{\underline{初始设置}：}$\Delta = \left\lfloor\dfrac{q}{t}\right\rfloor$，$\{S_i\}_{i=0}^{k-1} \xleftarrow{\$} \mathcal{R}_{\langle n, \textit{tern} \rangle}^k$

$ $

$ $

\textbf{\underline{加密输入}：}$M \in \mathcal{R}_{\langle n, t \rangle}$，$\{A_i\}_{i=0}^{k-1} \xleftarrow{\$} \mathcal{R}_{\langle n,q \rangle}^{k}$，$E \xleftarrow{\chi_\sigma} \mathcal{R}_{\langle n,q \rangle}$

$ $

\begin{enumerate}

\item 向上缩放 $M \longrightarrow \Delta M \text { } \in \mathcal{R}_{\langle n, q\rangle}$

\item 计算 $B = \sum\limits_{i=0}^{k-1}{(A_i \cdot S_i)} + \Delta  M + E \text{ } \in \mathcal{R}_{\langle n,q \rangle}$

\item $\textsf{GLWE}_{S,\sigma}(\Delta M + E) = (\{A_i\}_{i=0}^{k-1}, B) \text{ } \in \mathcal{R}_{\langle n,q \rangle}^{k + 1}$ 



\end{enumerate}

\end{tcolorbox}

\subsection{解密}
\label{subsec:glwe-dec}

\begin{enumerate}
\item 给定密文 $(\{A_i\}_{i=0}^{k-1}, B)$，其中 $B = \sum\limits_{i=0}^{k-1}{(A_i \cdot S_i)} + \Delta \cdot M + E \in \mathcal{R}_{\langle n,q \rangle}$，计算 $B - \sum\limits_{i=0}^{k-1}{(A_i \cdot S_i)} = \Delta \cdot M + E$。
\item 将多项式 $\Delta \cdot M + E \in \mathcal{R}_{\langle n,q \rangle}$ 的每个系数舍入到最接近的 $\Delta$ 的整数倍，即按基数 $\Delta$ 舍入，记为 $\lceil \Delta \cdot M + E \rfloor_{\Delta}$。这一操作可以成功消去 $E$，得到 $\Delta \cdot M$。需要注意的是，$E$ 的每个系数 $e_i$ 都必须满足 $|e_i| < \dfrac{\Delta}{2}$，才能在舍入时被消去。否则，$e_i$ 的部分高位会与明文 $m_i$ 系数的低位重叠，无法在解密时消去，从而破坏明文 $m_i$。
\item 计算 $\dfrac{\Delta \cdot M} {\Delta}$，等价于将 $\Delta \cdot M$ 的每个多项式系数按 $\Delta$ 向下缩放。
\end{enumerate}

$ $

GLWE 解密公式汇总如下：

$ $

\begin{tcolorbox}[title={\textbf{\tboxlabel{\ref*{subsec:glwe-dec}} GLWE 解密}}]
\textbf{\underline{解密输入}：}$\textsf{ct} = (\{A_i\}_{i=0}^{k-1}, B) \text{ } \in \mathcal{R}_{\langle n,q \rangle}^{k + 1}$

\begin{enumerate}
\item $\textsf{GLWE}^{-1}_{S,\sigma}(\textsf{ct}) = B - \sum\limits_{i=0}^{k-1}{(A_i \cdot S_i)} = \Delta  M + E \text{ } \in \mathcal{R}_{\langle n,q \rangle}$

\item 向下缩放
$\Bigl\lceil \dfrac{ \Delta  M + E }{\Delta}\Bigr\rfloor \bmod t = M \text{ } \in \mathcal{R}_{\langle n, t \rangle}$

\end{enumerate}

为了正确解密，噪声多项式 $E$ 的每个系数 $e_i$ 都应满足 $|e_i| < \dfrac{\Delta}{2}$。

\end{tcolorbox}


\subsubsection{讨论}

下面的大小比较仅描述表示方式与暴力搜索空间。具体安全性共同依赖于 $n,k,q$、误差和秘密分布及样本数量；单独增大某个维数并不构成安全归约。
\begin{enumerate}
\item \textbf{LWE} 是多项式环次数参数为 $n = 1$ 时的 GLWE 特例。此时，$\{A_i\}_{i=0}^{k-1}, \{S_i\}_{i=0}^{k-1}, E$ 和 $M$ 中的所有多项式都是零次常数多项式，而 $A_i$ 和 $S_i$ 分别含有 $k$ 个这样的常数，因此各自构成一个向量。
\item \textbf{RLWE} 是 $k = 1$ 时的 GLWE 特例。此时，秘密密钥 $S$ 是单个多项式 $S_0$，每次加密也只使用单个多项式 $A_0$ 作为公开量。
\item \textbf{$\bm{n}$ 的大小：}多项式 $S_i = s_{i,0} + s_{i,1}X + \cdots + s_{i, n-1}X^{n-1} $ 有 $n$ 个系数。若各系数独立均匀取三元值，则完整秘密有 $3^{nk}$ 种可能取值。这是穷举候选数量，并非最佳攻击的下界，也不是搜索假设（\autoref{subsec:lattice-overview}）的证明。更高次数的明文多项式可以表示更多系数，但实际吞吐量还取决于运算成本。
%\item \textbf{Size of $\bm{q}$:} Increasing $q$ increases the ciphertext molulus space, making it hard to attack the cryptosystem. 
%Also, a large modulo space creates an enough separation between the plaintext bits and the noise bits in the ciphertext, which reduces the growth rate of the noise bits across ciphertext multiplications (will be discussed in \autoref{subsubsec:glwe-mult-plain-discussion}) and thus we can do a more number of ciphertext multiplications before the noise bits overflow to the plaintext bit area. 
\item \textbf{$\bm{k}$ 的大小：}增大 $k$ 会增加秘密多项式 $(S_0, S_1, \ldots, S_{k-1})$ 和公开掩码多项式 $(A_0, A_1,\ldots, A_{k-1})$。每个密文仍只有一个消息多项式 $M$ 和一个误差多项式 $E$。比较两个 $n,k$ 选择时，需要考察其实际环或模结构假设与分布，而不能只数秘密系数。
\item \textbf{减小密文大小：}种子压缩密文可以保存新鲜独立的公开种子 $d$ 与 $B$。明确指定的公开展开函数 $G$ 生成全部 $kn$ 个掩码系数，并使用参数/索引域分离及无偏系数映射。其输入绝不是秘密。由于种子公开，普通隐藏种子的 PRG 安全性并不能说明 $(d,G(d))$ 可替换为带独立均匀 $A$ 的 $(d,A)$：检验 $A=G(d)$ 就能区分两者。对固定展开函数，所需联合假设比较的是 $(d,\langle G(d),S\rangle+E)$ 与均匀 $V$ 对应的 $(d,V)$，并须使用实际秘密/误差分布及样本数量。
\end{enumerate}

\noindent\textbf{明确的种子模型：}若展开是理想随机预言机 $H:\{0,1\}^h\to\mathcal R_{\langle n,q\rangle}^k$，经典攻击者作 $Q$ 次加密查询和 $Q_H$ 次预言机查询时，其猜测偏差满足
\[
\epsilon\le\delta(T',Q)+\min\left\{1,\frac{QQ_H+Q(Q-1)/2}{2^h}\right\}.
\]
这里 $\delta$ 是实际分布的判定 GLWE 优势，采用 \autoref{subsec:lattice-scheme} 的约定。模拟密文时独立选择种子，把挑战样本的均匀掩码编程到此前未赋值的 $H(d)$，再向身体加入选定消息。仅当种子此前已被查询，或重复了更早的加密种子时，模拟才失败。第 $j$ 次加密时至多有 $Q_H+j-1$ 个点已经占用，因此得到所列联合界。在该事件之外，真实样本产生真实加密，均匀样本身体则隐藏消息；失败时中止并作独立公平猜测。这证明了上述经典随机预言机结论，既没有不加假设地把具体 XOF 替换为理想预言机，也没有为量子预言机查询证明同一个界。关于公开 XOF 展开与秘密伪随机生成的区别，参见 \href{https://pq-crystals.org/kyber/data/kyber-specification-round3-20210804.pdf}{Kyber v3.02 说明第 4.3.1 与 5.4 节}。重复掩码还会暴露两个消息相位之差加上误差差值，因此种子的新鲜性至关重要。

\subsection{GLWE 的另一种形式}
\label{subsec:glwe-alternative}

下面给出 \tboxlabel{\ref*{subsec:glwe-enc}} 的另一种形式，在加密与解密公式中翻转每个 $A_iS_i$ 的符号：

\begin{tcolorbox}[title={\textbf{\tboxlabel{\ref*{subsec:glwe-alternative}} GLWE 密码体制的另一种形式}}]

\textbf{\underline{初始设置}：}$\Delta = \left\lfloor\dfrac{q}{t}\right\rfloor$，$\{S_i\}_{i=0}^{k-1} \xleftarrow{\$} \mathcal{R}_{\langle n, \textit{tern} \rangle}^k$

$ $

\par\noindent\rule{\textwidth}{0.4pt}


\textbf{\underline{加密输入}：}$M \in \mathcal{R}_{\langle n, t \rangle}$，$\{A_i\}_{i=0}^{k-1} \xleftarrow{\$} \mathcal{R}_{\langle n,q \rangle}^{k}$，$E \xleftarrow{\chi_\sigma} \mathcal{R}_{\langle n,q \rangle}$

$ $

\begin{enumerate}
\item 向上缩放 $M \longrightarrow \Delta M \text { } \in \mathcal{R}_{\langle n, q\rangle}$

\item 计算 $B = \textcolor{red}{-\sum\limits_{i=0}^{k-1}{(A_i \cdot S_i)}} + \Delta  M + E \text{ } \in \mathcal{R}_{\langle n,q \rangle}$

\item $\textsf{GLWE}_{S,\sigma}(\Delta M + E) = (\{A_i\}_{i=0}^{k-1}, B) \text{ } \in \mathcal{R}_{\langle n,q \rangle}^{k + 1}$ 

\end{enumerate}

\par\noindent\rule{\textwidth}{0.4pt}


\textbf{\underline{解密输入}：}$\textsf{ct} = (\{A_i\}_{i=0}^{k-1}, B) \text{ } \in \mathcal{R}_{\langle n,q \rangle}^{k + 1}$

\begin{enumerate}
\item $\textsf{GLWE}^{-1}_{S,\sigma}(\textsf{ct}) = B \text{ } \textcolor{red}{ + \sum\limits_{i=0}^{k-1}{(A_i \cdot S_i)}} = \Delta  M + E \text{ } \in \mathcal{R}_{\langle n,q \rangle}$

\item 向下缩放
$\Bigl\lceil \dfrac{ \Delta  M + E }{\Delta}\Bigr\rfloor \bmod t = M  \text{ } \in \mathcal{R}_{\langle n, t \rangle}$

\end{enumerate}

为了正确解密，噪声多项式 $E$ 的每个系数 $e_i$ 都应满足 $|e_i| < \dfrac{\Delta}{2}$。

\end{tcolorbox}

即使 $A_iS_i$ 项的符号发生翻转，解密阶段也会通过加上相反符号的对应项将其抵消，因此仍保持与原形式相同的解密正确性。





\subsection{公钥加密}
\label{subsec:glwe-public-key-enc}

\textbf{适用范围：}本小节限定 $k=1$，即标准 RLWE 公钥结构，并要求足够大的二幂环维数 $n$，以及能够论证所需判定 RLWE 假设的参数和分布族。下面保留只有一个分量的向量记号。这不是适用于任意 GLWE 维数或 $n=1$ 的 LWE 特例的通用公钥构造；这些推广需要合适的公开矩阵构造及各自的安全论证。与前面的设置相同，展示的解密界假设 $t\mid q$。


\autoref{subsec:glwe-enc} 中的加密方案假设，每个明文都由秘密密钥持有者加密。本节介绍公钥加密方案：生成与秘密密钥相对应的公钥后，任何知道公钥的人都可以加密明文，而只有秘密密钥持有者能够解密。其基本思路是在初始化阶段预先计算部分加密所需的分量，并作为公钥公开。在实际加密阶段，将公钥乘以额外的随机量 $U$，再加入额外噪声 $E_1, \vec{E}_2$，为每个密文产生不可预测的随机性。具体方案如下：

\begin{tcolorbox}[title={\textbf{\tboxlabel{\ref*{subsec:glwe-public-key-enc}} RLWE 公钥加密（$k=1$）}}]

\textbf{\underline{初始设置}：}
\begin{itemize}
\item 尺度因子 $\Delta = \left\lfloor\dfrac{q}{t}\right\rfloor$
\item 秘密密钥 $\vec{S} = \{S_i\}_{i=0}^{k-1} \xleftarrow{\$} \mathcal{R}_{\langle n, \textit{tern} \rangle}^k$
\item 按如下方式生成公钥对 $(\mathit{PK}_1, \vv{\mathit{PK}}_2) \in \mathcal{R}_{\langle n, q \rangle}^{k+1}$：

$\vec{A} = \{A_i\}_{i=0}^{k-1} \xleftarrow{\$} \mathcal{R}_{\langle n, q \rangle}^k$, \text{ } $E \xleftarrow{\chi_{\sigma}} \mathcal{R}_{\langle n, q \rangle}$

$\mathit{PK}_1 = \vec{A} \cdot \vec{S} + E \in \mathcal{R}_{\langle n, q \rangle}$

$\vv{\mathit{PK}}_2 = \vec{A} \in \mathcal{R}_{\langle n, q \rangle}^{k}$

\end{itemize}

\par\noindent\rule{\textwidth}{0.4pt}

\textbf{\underline{加密输入}：}$M \in \mathcal{R}_{\langle n, t \rangle}$，\text{ }$U \xleftarrow{\$} \mathcal{R}_{\langle n,\textit{tern} \rangle}, \text{ } E_1 \xleftarrow{\chi_{\sigma}} \mathcal{R}_{\langle n,q \rangle}, \text{ } \vec{E}_2 \xleftarrow{\chi_{\sigma}} \mathcal{R}_{\langle n,q \rangle}^k$

$ $

\begin{enumerate}

\item 向上缩放 $M \longrightarrow \Delta M \text { } \in \mathcal{R}_{\langle n, q\rangle}$

\item 计算：

$B = \mathit{PK}_1\cdot U + \Delta  M + E_1 \text{ } \in \mathcal{R}_{\langle n,q \rangle}$

$\vec{D} = \vv{{\mathit{PK}}}_2 \cdot U + \vec{E}_2 \in \mathcal{R}_{\langle n,q \rangle}^{k}$ \textcolor{red}{ $\rhd$ $\vv{\mathit{PK}}_2 \cdot U$ 将 $\vv{\mathit{PK}}_2$ 的每个元素乘以 $U$}

\item $\textsf{GLWE}_{S,\sigma}(\Delta M + E_{\mathit{all}}) = (\vec{D}, B) \text{ } \in \mathcal{R}_{\langle n,q \rangle}^{k+1}$ \textcolor{red}{ $\rhd$ 其中 $E_{\mathit{all}} = E \cdot U + E_1 - \vec{E}_2\cdot\vec{S}$}



\end{enumerate}


\par\noindent\rule{\textwidth}{0.4pt}

\textbf{\underline{解密输入}：}$\textsf{ct} = (\vec{D}, B) \text{ } \in \mathcal{R}_{\langle n,q \rangle}^{k + 1}$

\begin{enumerate}
\item $\textsf{GLWE}^{-1}_{S,\sigma}(\textsf{ct}) = B - \vec{D} \cdot \vec{S} = \Delta  M + E_{\mathit{all}} \text{ } \in \mathcal{R}_{\langle n,q \rangle}$

\item 向下缩放
$\Bigl\lceil \dfrac{ \Delta  M + E_{\mathit{all}} }{\Delta}\Bigr\rfloor \bmod t = M  \text{ } \in \mathcal{R}_{\langle n, t \rangle}$

\end{enumerate}

为了正确解密，噪声多项式 $E_{\mathit{all}}$ 的每个系数 $e_i$ 都应满足 $|e_i| < \dfrac{\Delta}{2}$。

\end{tcolorbox}


解密过程第一步中的等式推导如下：

$\textsf{GLWE}^{-1}_{S,\sigma}\bm{(}\text{ } \textsf{ct} = (\vec{D}, B) \text{ }\bm{)} = B - \vec{D}\cdot\vec{S}$

$ = (\mathit{PK}_1\cdot U + \Delta  M + E_1) - (\vv{\mathit{PK}}_2 \cdot U + \vec{E}_2)\cdot \vec{S} $

$=  (\vec{A} \cdot \vec{S} + E)\cdot U + \Delta M + E_1 - (\vec{A}\cdot U)\cdot \vec{S} - \vec{E}_2\cdot\vec{S}$

$=  (U\cdot\vec{A}) \cdot \vec{S} + E \cdot U + \Delta M + E_1 - (U\cdot \vec{A}) \cdot \vec{S} - \vec{E}_2\cdot\vec{S}$

$= \Delta M + E \cdot U + E_1 - \vec{E}_2\cdot\vec{S}$

$= \Delta M + E_{\mathit{all}}$ \textcolor{red}{ $\rhd$ 其中 $E_{\mathit{all}} = E \cdot U + E_1 - \vec{E}_2\cdot\vec{S}$}


$ $

\para{安全性：}在 $k=1$ 时，将公钥写成 $(A,P)$，其中 $P=AS+E$；密文为 $(D,B)=(AU+E_2,PU+E_1+\Delta M)$。假设针对实际 $S$ 分布及误差 $E$ 的判定 RLWE 成立，也假设针对实际 $U$ 分布的两样本判定 RLWE 成立，其两个误差分别服从 $E_2,E_1$ 的分布。本小节中两个秘密均为三元分布，三种误差分布相同，因此同一个相应的两样本 RLWE 假设可覆盖两次使用。仅有针对高斯秘密的定理并不能推出这一点。

单挑战 IND-CPA 证明使用两个混合。首先，将 $P=AS+E$ 替换为独立均匀环元素。归约将一条 RLWE 样本嵌入公钥，自行选择新鲜 $U,E_1,E_2$ 就能执行全部加密，不需要知道 $S$。随后，在均匀 $(A,P)$ 下，将挑战的 $(D,B-\Delta M)$ 替换为独立均匀环元素。归约取得共享同一个未知秘密 $U$ 的两条样本 $(A,Z_2),(P,Z_1)$，以其掩码作为公钥，等挑战消息选定后返回 $(Z_2,Z_1+\Delta M)$。真实样本给出 $Z_2=AU+E_2$ 和 $Z_1=PU+E_1$；均匀样本则完全隐藏选中的消息。其他加密都可用新鲜、已知的随机性公开执行。因此，采用 \autoref{subsec:lattice-scheme} 的猜测偏差约定，有
\[
\epsilon_{\mathrm{PKE}}\le\delta_S(T'_1,1)+\delta_U(T'_2,2).
\]
这两项使用所述实际分布，并计入归约开销。这是限定构造的完整条件性 CPA 归约，并非某个数值参数族的认证，也不涵盖依赖秘密的公钥加密查询。参见 \href{https://web.eecs.umich.edu/~cpeikert/pubs/ideal-lwe.pdf}{Lyubashevsky--Peikert--Regev，第 1.1 节 Applications}。对于未加限制的向量版本，第二个混合需要的是以 $U$ 为秘密、具有 $k+1$ 条样本的秩一 RLWE，而不仅是第一个混合所用的秩 $k$ GLWE 假设。

抵消运算另外给出了正确性恒等式，其中 $E_{\mathit{all}}=EU+E_1-E_2S$。

在负循环环中，对提升后的系数有
\[
\|E_{\mathit{all}}\|_\infty\le\|E\|_\infty\|U\|_1+\|E_1\|_\infty+\|E_2\|_\infty\|S\|_1.
\]
当 $t\mid q$ 时，右端小于 $\Delta/2$ 即足以保证解密。

如果不限制适用范围，那么在 $n=1$ 且 $U$ 为三元数时，即使 LWE 向量维数 $k$ 很大，这个构造也不安全：攻击者只需尝试 $U\in\{-1,0,1\}$，检验所有 $D_i-A_iU$ 是否很小，再解码 $B-\mathit{PK}_1U$。特别地，$U=0$ 以 $1/3$ 的概率发生，此时直接得到 $B=\Delta m+E_1$。因此，不能从展示的代数公式推出一般 GLWE 公钥安全性。
